\paragraph{Fourier--\(p\)-adic spectra of elliptic projection gates.}
We study the \(m\)-fold covering of a confocal ellipse obtained from the unit circle by the map \(e^{i\theta}\mapsto e^{im\theta}+c e^{-im\theta}\). Its logarithmic speed has an exceptionally sparse Fourier expansion, and this sparseness turns the rational parameter \(c\) into a package of local valuation invariants. The same spectral data also determine the minimal size of any finite auxiliary register that separates the \(m\) sheets of the covering, while finite congruence tests necessarily miss sufficiently deep \(p\)-adic perturbations.

\begin{definition}[Elliptic projection gate and logarithmic speed spectrum]\label{def:group-jg-elliptic-gate-logspeed-spectrum}
Let \(c\in(0,1)\cap\QQ\) and \(m\in\NN\). Define
\[
\gamma_{c,m}(\theta):=J_{c,m}(e^{i\theta})=e^{im\theta}+c\,e^{-im\theta}
\qquad (\theta\in\RR/2\pi\ZZ).
\]
Its image is the ellipse
\[
E_c:=\bigl\{(1+c)\cos t+i(1-c)\sin t:\ t\in\RR/2\pi\ZZ\bigr\}\subset\CC.
\]
The logarithmic speed defect is
\[
u_{c,m}(\theta):=\log\abs{\gamma_{c,m}'(\theta)}-\log m
=\log\abs{1-c\,e^{-2im\theta}},
\]
and its Fourier coefficients are
\[
\widehat u_{c,m}(k):=\frac{1}{2\pi}\int_{0}^{2\pi}u_{c,m}(\theta)e^{-ik\theta}\,\dd\theta
\qquad (k\in\ZZ).
\]
\end{definition}

\begin{theorem}[Prime-spectral rigidity of elliptic projection gates]\label{thm:group-jg-elliptic-gate-prime-spectrum-rigidity}
In the setting of Definition \ref{def:group-jg-elliptic-gate-logspeed-spectrum} one has:
\begin{enumerate}
  \item \textbf{Sparse support.}
  The zero Fourier mode vanishes, \(\widehat u_{c,m}(0)=0\), and for every \(n\ge 1\),
  \[
  \widehat u_{c,m}(\pm 2mn)=-\frac{c^n}{2n}.
  \]
  In particular, \(\widehat u_{c,m}(k)=0\) whenever \(k\notin 2m\ZZ\), and the nonzero support is exactly \(2m\ZZ\setminus\{0\}\).

  \item \textbf{\(p\)-adic slopes.}
  For every prime \(p\) and every \(n\ge 1\),
  \[
  \nu_p\!\bigl(\widehat u_{c,m}(2mn)\bigr)
  =n\,\nu_p(c)-\nu_p(n)-\nu_p(2).
  \]
  In particular, for \(p\neq 2\),
  \[
  \nu_p\!\bigl(\widehat u_{c,m}(2mn)\bigr)=n\,\nu_p(c)-\nu_p(n).
  \]

  \item \textbf{Parameter reconstruction.}
  The full Fourier family uniquely determines \(m\) and \(c\). More precisely,
  \[
  m=\frac12\gcd\{\,|k|:\ \widehat u_{c,m}(k)\neq 0\,\},
  \qquad
  c=-2\,\widehat u_{c,m}(2m).
  \]
  For each odd prime \(p\),
  \[
  \nu_p(c)
  =
  \lim_{\substack{n\to\infty\\ p\nmid n}}
  \frac{\nu_p\!\bigl(\widehat u_{c,m}(2mn)\bigr)}{n}.
  \]
\end{enumerate}
Consequently, \((c,m)\), equivalently the valuation data \(\bigl((\nu_p(c))_p,(\nu_p(m))_p\bigr)\), is a complete invariant of the logarithmic speed spectrum.
\end{theorem}

\begin{proof}
\textbf{(1)} Since \(0<c<1\), the identity
\[
\log(1-z)=-\sum_{n\ge 1}\frac{z^n}{n}
\]
holds for \(|z|<1\). Substituting \(z=c e^{-2im\theta}\) and taking real parts gives
\[
u_{c,m}(\theta)=\Re\log(1-c e^{-2im\theta})
=-\sum_{n\ge 1}\frac{c^n}{n}\cos(2mn\theta).
\]
Using
\[
\cos(2mn\theta)=\frac12\bigl(e^{i2mn\theta}+e^{-i2mn\theta}\bigr)
\]
we obtain
\[
\widehat u_{c,m}(\pm 2mn)=-\frac{c^n}{2n}
\qquad (n\ge 1),
\]
while every other Fourier coefficient, including the zeroth one, vanishes. This proves the support claim.

\textbf{(2)} By part (1),
\[
\widehat u_{c,m}(2mn)=-\frac{c^n}{2n}.
\]
Writing \(c=a/b\) in lowest terms yields
\[
\widehat u_{c,m}(2mn)=-\frac{a^n}{2b^n n},
\]
hence
\[
\nu_p\!\bigl(\widehat u_{c,m}(2mn)\bigr)
=n\bigl(\nu_p(a)-\nu_p(b)\bigr)-\nu_p(n)-\nu_p(2)
=n\,\nu_p(c)-\nu_p(n)-\nu_p(2).
\]

\textbf{(3)} By part (1), the nonzero frequencies are exactly \(\{\pm 2m,\pm 4m,\ldots\}\), so their greatest common divisor is \(2m\), which recovers \(m\). The first nonzero harmonic gives
\[
\widehat u_{c,m}(2m)=-\frac{c}{2},
\]
hence \(c=-2\,\widehat u_{c,m}(2m)\). The odd-prime slope formula is the specialization of part (2) along integers \(n\) with \(p\nmid n\), for which \(\nu_p(n)=\nu_p(2)=0\). Therefore the Fourier family determines \(m\), \(c\), and hence the full valuation vectors of both \(c\) and \(m\).
\end{proof}

\begin{theorem}[Outer-function rigidity of the logarithmic speed profile]\label{thm:group-jg-outer-function-rigidity-joukowsky-ellipse}
Let \(\Phi\) be holomorphic on the exterior disk \(\{\,|z|>1\,\}\) and meromorphic at infinity with Laurent expansion
\[
\Phi(z)=A z^m+O(z^{m-1})
\qquad (z\to\infty),
\]
where \(m\in\NN\) and \(A\in\CC^\times\). Assume that the boundary trace
\[
\gamma(\theta):=\Phi(e^{i\theta})
\]
is a real-analytic immersion and that, for some \(c\in(0,1)\),
\[
\log\abs{\gamma'(\theta)}
=\log(m\abs{A})+\log\abs{1-c e^{-2im\theta}}
\qquad (\theta\in\RR/2\pi\ZZ).
\]
Then there exists \(B\in\CC\) such that
\[
\Phi(z)=A\bigl(z^m+c z^{-m}\bigr)+B.
\]
Therefore \(\gamma\) is, up to translation, rotation, and dilation, the standard \(m\)-fold parametrization of the ellipse \(E_c\).
\end{theorem}

\begin{proof}
Define
\[
F(z):=\frac{\Phi'(z)}{mA z^{m-1}}.
\]
Because \(\Phi(z)=A z^m+O(z^{m-1})\) at infinity, \(F\) is holomorphic on \(\{|z|>1\}\) and satisfies \(F(\infty)=1\). For \(z=e^{i\theta}\) we have
\[
\gamma'(\theta)=\ii e^{i\theta}\Phi'(e^{i\theta}),
\]
so \(\abs{\gamma'(\theta)}=\abs{\Phi'(e^{i\theta})}\). The hypothesis therefore yields
\[
\abs{F(e^{i\theta})}=\abs{1-c e^{-2im\theta}}=\abs{1-c z^{-2m}}.
\]
Since \(c<1\), the factor \(1-c z^{-2m}\) has no zeros on \(|z|\ge 1\). Hence
\[
G(z):=\frac{F(z)}{1-c z^{-2m}}
\]
is holomorphic on \(\{|z|>1\}\), continuous up to \(|z|=1\), and satisfies
\[
\abs{G(e^{i\theta})}=1,
\qquad
G(\infty)=1.
\]
Passing to the unit disk by \(H(\zeta):=G(1/\zeta)\), we obtain a holomorphic function on \(\DD\) with \(H(0)=1\) and \(\abs{H(\zeta)}=1\) on \(|\zeta|=1\). By the maximum modulus principle, \(H\equiv 1\), hence \(G\equiv 1\). Therefore
\[
F(z)=1-c z^{-2m},
\qquad\text{so}\qquad
\Phi'(z)=mA\bigl(z^{m-1}-c z^{-m-1}\bigr).
\]
Integrating gives
\[
\Phi(z)=A\bigl(z^m+c z^{-m}\bigr)+B
\]
for some \(B\in\CC\), as claimed.
\end{proof}

\begin{proposition}[Minimal auxiliary registers for \(m\)-sheeted elliptic projections]\label{prop:group-jg-elliptic-gate-minimal-register}
Let \(c\in(0,1)\), \(m\in\NN\), and consider the covering \(\gamma_{c,m}:\TT\to E_c\). Suppose that \(R\) is a finite set and that there exists an injective map
\[
\iota:\TT\longrightarrow E_c\times R
\]
such that \(\pi_{E_c}\circ\iota=\gamma_{c,m}\). Then:
\begin{enumerate}
  \item \(\abs{R}\ge m\). In particular, any finite auxiliary register that separates the \(m\) sheets carries at least \(\log m\) units of information.

  \item Suppose in addition that \(\iota\) is equivariant for the deck group \(\mu_m:=\{\zeta\in\CC:\zeta^m=1\}\cong\ZZ/m\ZZ\), in the sense that there is a permutation representation \(\rho:\mu_m\to\Sym(R)\) with
  \[
  \iota(\zeta z)=\bigl(\gamma_{c,m}(z),\,\rho(\zeta)(r(z))\bigr).
  \]
  Then \(\rho\) contains an orbit of length \(m\). If, moreover, \(R\) is endowed with a finite abelian group structure and \(\rho\) is realized by translations, then \(R\) contains a cyclic subgroup of order \(m\), so that
  \[
  \ZZ/m\ZZ\cong \prod_{p}\ZZ/p^{\nu_p(m)}\ZZ,
  \qquad
  \log m=\sum_p \nu_p(m)\log p.
  \]
\end{enumerate}
\end{proposition}

\begin{proof}
For every \(w\in E_c\), the fiber \(\gamma_{c,m}^{-1}(w)\) contains exactly \(m\) points, since the parametrization of \(E_c\) by \(t\mapsto (1+c)\cos t+i(1-c)\sin t\) is one-to-one modulo \(2\pi\), and \(\theta\mapsto m\theta\) is an \(m\)-fold covering of the circle. If \(\iota\) is injective, then the \(m\) points above \(w\) have the same first coordinate in \(E_c\times R\), so their second coordinates must be distinct. Hence \(\abs{R}\ge m\).

For the equivariant statement, the deck group \(\mu_m\) acts freely and transitively on each fiber of \(\gamma_{c,m}\). Therefore the induced action on labels must be free and transitive on the corresponding subset of \(R\), which is equivalent to the existence of an orbit of length \(m\). If the action is implemented by translations in a finite abelian group \(R\), such an orbit is a coset of a cyclic subgroup of order \(m\). The displayed decomposition is the standard primary decomposition of a finite cyclic group, and the logarithmic identity follows by taking logarithms of the factorization of \(m\).
\end{proof}

\begin{corollary}[Prime-spectral classification of rational elliptic projection gates]\label{cor:group-jg-elliptic-gate-prime-spectrum-classification}
In the equivariant-register setting of Proposition \ref{prop:group-jg-elliptic-gate-minimal-register}, a rational elliptic projection gate \(J_{c,m}\) with \(c\in(0,1)\cap\QQ\) and \(m\in\NN\) is determined, up to rotation, translation, and relabeling of the auxiliary register, by the valuation data
\[
\Bigl((\nu_p(c))_{p\ \mathrm{prime}},\,(\nu_p(m))_{p\ \mathrm{prime}}\Bigr).
\]
Conversely, any finitely supported valuation vectors \((\alpha_p)_p,(\beta_p)_p\in\bigoplus_p\ZZ\) with
\[
c=\prod_p p^{\alpha_p}\in(0,1)\cap\QQ,
\qquad
m=\prod_p p^{\beta_p}\in\NN
\]
determine a unique such isomorphism class.
\end{corollary}

\begin{proof}
Theorem \ref{thm:group-jg-elliptic-gate-prime-spectrum-rigidity} recovers \(m\) and \(c\) from the Fourier family, hence it recovers the full valuation vector \((\nu_p(c))_p\). Proposition \ref{prop:group-jg-elliptic-gate-minimal-register} shows that any equivariant auxiliary register must contain a cyclic orbit of size \(m\), so the valuation data of \(m\) are intrinsic as well. Conversely, finitely supported valuation vectors uniquely determine the rational number \(c\) and the integer \(m\), and therefore determine the corresponding gate.
\end{proof}

\begin{theorem}[\(p\)-adic blindspots for finite congruence audits]\label{thm:group-jg-elliptic-gate-padics-blindspot}
Fix an odd prime \(p\ge 3\), an integer \(K\ge 1\), and a truncation level \(N\ge 1\). Then there exist infinitely many pairs of distinct rational numbers
\[
c,\ c'\in (0,1)\cap\QQ\cap\ZZ_{(p)}^\times
\]
such that
\[
\widehat u_{c,1}(2n)\equiv \widehat u_{c',1}(2n)\pmod{p^K\ZZ_{(p)}}
\qquad (1\le n\le N).
\]
Thus no audit that inspects only finitely many Fourier coefficients modulo \(p^K\) can exclude all sufficiently deep \(p\)-adic perturbations of the parameter.
\end{theorem}

\begin{proof}
Let
\[
M:=K+\max_{1\le n\le N}\nu_p(n).
\]
Choose any integer \(b>1+p^M\) with \(p\nmid b\), and set
\[
c:=\frac{1}{b},
\qquad
c':=\frac{1+p^M}{b}.
\]
Then \(c\neq c'\), both numbers lie in \((0,1)\cap\QQ\), and both are \(p\)-adic units in \(\ZZ_{(p)}\). Moreover,
\[
c'-c=\frac{p^M}{b}\in p^M\ZZ_{(p)}.
\]
For \(n\ge 1\), the binomial expansion gives
\[
c'^n-c^n=\sum_{j=1}^{n}\binom{n}{j}c^{n-j}\left(\frac{p^M}{b}\right)^j.
\]
The \(j=1\) term has \(p\)-adic valuation \(M+\nu_p(n)\), while every term with \(j\ge 2\) has valuation at least \(2M\). Since \(M\ge \nu_p(n)\) for \(1\le n\le N\), every summand has valuation at least \(M+\nu_p(n)\), and therefore
\[
\nu_p(c'^n-c^n)\ge M+\nu_p(n)\ge K+\nu_p(n)
\qquad (1\le n\le N).
\]
By Theorem \ref{thm:group-jg-elliptic-gate-prime-spectrum-rigidity}(1) with \(m=1\),
\[
\widehat u_{c,1}(2n)=-\frac{c^n}{2n},
\qquad
\widehat u_{c',1}(2n)-\widehat u_{c,1}(2n)
=-\frac{c'^n-c^n}{2n}.
\]
Because \(p\ge 3\), one has \(\nu_p(2)=0\), so for \(1\le n\le N\),
\[
\nu_p\!\bigl(\widehat u_{c',1}(2n)-\widehat u_{c,1}(2n)\bigr)
\ge K.
\]
This proves the congruence claim. There are infinitely many admissible integers \(b\), hence infinitely many such pairs \((c,c')\).
\end{proof}

\endinput
